
% language=us runpath=texruns:manuals/mwe

\usemodule[scite]

\setuplayout
  [tight]

% \showframe

% \setuplayout
%   [middle]

\setuphead
  [section]
  [style=\bfa]

\setupbodyfont
  [pagella,12pt]

\setupwhitespace
  [big]

\setupalign
  [mathbookpasses]

\setupdocument
  [title=No Title,
   before=\setups{titlepage}]

\startsetups[titlepage]
\startMPpage
    picture p ; p := image (
        draw textext("\tt\bf\darkyellow M{\darkred N}WE")   xysized (PaperWidth,PaperHeight/3) shifted (0,3PaperHeight/3);
        draw textext("\tt\bf\darkyellow {\darkgreen N}WME") xysized (PaperWidth,PaperHeight/3) shifted (0,2PaperHeight/3);
        draw textext("\tt\bf\darkyellow ME{\darkblue N}W")  xysized (PaperWidth,PaperHeight/3) shifted (0,1PaperHeight/3);
    ) ;
    fill boundingbox p withcolor .8 white;
    draw p xysized (PaperWidth-10mm,PaperHeight-10mm) shifted (0,5mm) ;
\stopMPpage
\stopsetups

\noheaderandfooterlines

\startdocument

\startsubject[title=A test file]

This is how a minimal (non) working example looks like. The more minimal it is,
the more likely you will get an answer. Remove everything from your examples
which isn't necessary to show the problem:

\startitemize[packed]
    \startitem the page layout \stopitem
    \startitem the font settings \stopitem
    \startitem the title page \stopitem
    \startitem settings than don't matter \stopitem
\stopitemize

Just a simple:

\starttyping[option=TEX]
\starttext
    a dozen lines at most
\stoptext
\stoptyping

Like this:

\starttyping[option=TEX]
\starttext
    an example
\stoptext
\stoptyping

and look at it a few times to see what can be left out. If you have a complex
workflow, and more files are needed, make a zip and put it somewhere to download;
don't spam the list with files (reports and explanations are fine).

\stopsubject

\startsubject[title=An explanation]

Because \TEX\ is also about look and feel it can be that your expectations
don't match what \CONTEXT\ generates. In that case explain what you needed
and what you expect. Maybe make a drawing or fake output. There is no need
to come up witt some interface, if needed we will see how and
where a possible feature (option) fits in best.

When you have specific wishes or are in a hurry, also be prepared to test as soon
as a patch or update is provided. Don't expect us to stay tuned to your problem
for too long.

\stopsubject

\startsubject[title=No slop]

We see no problem in users asking some large language model driven
tool. But don't spam us with its outcome. We start from a problem not a
proposed solution. Of course fixing a simple typo or wrong variable is
no problem but the internals of \CONTEXT\ let alone the engine
mix badly with the kind of reasoning or lack of understanding per 2026.

Also, make clear that if you have a suggestion that it comes from such an agent.
If you don't really understand what is going on, just say so and don't rely on an
language model to make you wiser with hundred percent accuracy.

\stopsubject

\stopsubject

Hans (also for Wolfgang, Mikael, \unknown)

\stopdocument
